\section{Vom Photon zum Kompensator}
\label{sec:question}

Angefangen hat das Ganze für mich beim Photon.

Was muss eine Konstruktion beschreiben können, damit sie als Modell für Licht infrage kommt? Unser erster Prüfstein ist eine ebene elektromagnetische Welle im Vakuum. In Einheiten mit \(c=1\) sind ihre Feldvektoren gleich lang und orthogonal. An diesen beiden Eigenschaften prüfen wir die geplante Operation. Die Frage nach einem quantisierten Lichtzustand schließt sich an diesen klassischen Vergleich an.

Im Vorwort stand die Frage, was ein Heisenberg-Kompensator für die Physik bedeuten könnte. Der erste Schritt ist kleiner. Wir suchen eine Operation, die einen bestimmten Anteil eines Zustands aufhebt und den Rest erhält. Wie kann sie die beiden Anteile unterscheiden?
